The channel-selection analysis indicated that best-channel selection concentrated on frontal and frontopolar sites in both corpora (Table~\ref{tab:channel_selection}). This pattern was most directly interpretable for Cao2018, where the full recorded montage was available and selection genuinely concentrated on the frontopolar channels FP1 (52.6\%) and FP2 (26.3\%), making frontal dominance a meaningful empirical result rather than a consequence of limited channel availability. In contrast, the Raja results required more cautious interpretation because only frontal channels (including E3, E9, and E22) were offered as candidate selections, so frontal confinement in that corpus was partly by construction. The two proposed variants also selected the same channel in 89.1\% of Raja sessions and 82.8\% of Cao2018 sessions, far more often than other method pairs, indicating that the proposed thresholding strategy produced relatively stable channel choices. Together, these findings imply that a small frontal or frontopolar montage is sufficient for threshold-based blink detection, which simplifies hardware requirements and experimental setup without requiring a dense whole-scalp montage.
